% 第 17 章 Memory and Mixed Signal Design Analysis（原书 17-508）
\bichapter{存储器与混合信号设计分析}{Memory and Mixed Signal Design Analysis}
\label{chap:mms}

\section{引言}
\subsection*{Introduction}

在晶体管级执行的存储器与混合信号（Memory and Mixed Signal，MMX）分析，在当今先进半导体设计场景中日益重要。可从精确电源完整性分析获益的主要 MMX 应用包括：
\begin{itemize}
  \item 高性能 I/O，例如 DDR2、PCI-X、USB、Serdes、PLL 与 DLL
  \item 嵌入式存储器宏，例如 compiled memory、寄存器文件、TLB 与 cache memory
  \item 存储器芯片，例如 CAM、DRAM、SRAM 与 Flash memory
\end{itemize}

MMX 建模的优势包括：
\begin{itemize}
  \item 在混合了定制宏与基于单元数字模块的 SoC 上进行全芯片动态分析
  \item 大容量与较短运行时间（相较快速 SPICE 方法）
  \item 基于版图的 GUI，便于调试与诊断
\end{itemize}

有关在存储器与定制宏上以晶体管级精度执行精确 MMX 功耗分析、电源/地网格检查、静态 IR drop 分析，以及电迁移与动态压降分析的详细信息，请参阅 Apache \textit{Totem User Manual}。采用该详细表征建模方法时，器件模型在空间上（相对位置）与时间上（相对开关时刻）均保持准确。

有关在 RedHawk 中创建物理 GDS 模型的信息，请参阅附录 E 中 \cmd{gds2rh}/\cmd{gds2def} 一节。
